\chapter{Что модель умеет и чего не умеет}\label{ch:results}

\section{Как модель росла}

Модель собиралась шагами, и каждый шаг чему-то научил. Коротко:
\begin{enumerate}
\item \textbf{Последовательный автомат.} Пластинки появляются по одной, место выбирается с весом
  $e^{\beta g}$, поле соседей упругое, обрезанное потолком. Научил: поле пластинок само собирает
  окружные пакеты и радиальные «колоды»; направление цепочки переключается при наклоне пластинок
  $45^\circ$. Не научил: порога нет --- кривая RHF$(\sigma)$ пологая.
\item \textbf{Пластический вклад нагрузки.} Решатель на Фурье показал, что пластичность втрое
  усиливает выбор ориентации ($\Delta g\approx1.0\,\sigma_\theta$). Порога всё равно нет: зависимость линейна.
\item \textbf{Кинетика.} Время, остывание, диффузия водорода, теория зарождения относительно TSSP.
  Выпадение само идёт вдоль TSSP, баланс водорода сходится.
\item \textbf{Измеренный сдвиг и фора зёрен.} $0.08\,^\circ$C/МПа из опыта и фора окружного
  семейства. Появился резкий порог «фора/0.08» --- как ступенька Cinbiz.
\item \textbf{Границы зёрен} с форой из теории зарождения и \textbf{рост во времени} с переходом
  через границы --- доля межзёренных как у Son без подгонки.
\item \textbf{Свободная ось и ореолы} (эта сессия) --- пластичность вместо потолка; появились
  ступенчатые стопки.
\end{enumerate}

\section{Сводка проверок}

\begin{table}[htbp]\centering\small
\caption{Модель (кинетический автомат, до ореолов) против опыта. Подгонялись только $B$, $\Delta_0$
и фора зёрен Zircaloy-4; остальное --- проверки.}\label{tab:checks}
\begin{tabularx}{\linewidth}{>{\raggedright}p{4.0cm}>{\raggedright}p{3.5cm}>{\raggedright}X c}\toprule
закономерность & опыт & модель & итог\\\midrule
ступенька RHF у Cinbiz & 0 до 145~МПа, $\sim$1 выше 177 & 0.02 при 150, 0.68--0.85 при 175 & подогнано\\
Lepine: RHF при 0/200/250~МПа & 0.07 / 0.57 / 0.77 & 0.02 / 0.72 / 0.66 & средне\\
порог от водорода (Desquines) & 161--175~МПа при 100--450 ppm & 174--177~МПа & \textcolor{hydr}{да}\\
мало водорода без нагрузки & RHF 0.24--0.30 при 20--40 ppm (Son) & 0.26 при 60 ppm & \textcolor{hydr}{да}\\
неполное растворение & RHF растёт с $T_{\max}$ & 0.30 $\to$ 0.65 (350 $\to$ 420\,$^\circ$C) & \textcolor{hydr}{да}\\
доля межзёренных (Son) & 0.92 CWSR; 0.62 у радиальных & 0.91; 0.60--0.65 & \textcolor{hydr}{да}\\
начало переориентации Zr--Nb & $\approx80$~МПа (Son) & $\approx81$~МПа с измеренной форой 5\,$^\circ$C & \textcolor{hydr}{да}\\
ширина перехода (Son) & $>60$~МПа & 15--25~МПа & \textcolor{rad}{нет}\\
скорость охлаждения (Min) & медленнее --- больше радиальных & в пределах разброса & \textcolor{rad}{нет}\\
двухосность (Cinbiz) & порог 155 $\to$ 75~МПа & не меняется (2D и 3D) & \textcolor{rad}{нет}\\
стопки из мелких пластинок & Puls, гл.~3 & отдельные пластинки & \textcolor{rad}{нет}\\
\bottomrule
\end{tabularx}
\end{table}

\begin{idea}
Модель хорошо описывает то, что решает \emph{кто начал первым}: порог и его независимость от
водорода, неполное растворение, долю межзёренных, начало переориентации. Хуже --- то, что решает
\emph{как гидрид растёт и собирается}: ширину перехода, влияние скорости охлаждения,
морфологию стопок.
\end{idea}

\begin{exercise}
Почему фора в градусах даёт порог, не зависящий от водорода? (Подсказка: от водорода зависит
\emph{температура}, при которой начнётся выпадение, но не то, какое семейство начнёт первым.)
Почему тогда при 60 ppm радиальных много без нагрузки?
\end{exercise}

\section{Первые результаты с ореолами}

\fig{F23_Structures}{Структуры автомата при 0, 125, 175 и 250~МПа. Сверху --- старая модель
(потолок $\pm180$~МПа): радиальность даётся отдельными крутыми пластинками, рассыпанными по полю.
Снизу --- с пластическими ореолами: появились длинные радиальные стопки из мелких пластинок разной
ориентации. Параметры старые, не перекалиброванные.}

С ореолами при тех же параметрах (рис.~\ref{fig:F23_Structures}):
\begin{itemize}
\item \textbf{Морфология.} При 175--250~МПа радиальные гидриды --- стопки длиной до $\sim50$~мкм,
  сложенные из пластинок, из которых круче $45^\circ$ лишь половина. Радиальной получается
  стопка целиком: следующая пластинка садится у конца предыдущей, со сдвигом. Так описывает
  крупные гидриды Пульс: «стопки смещённых пластинок», видимая плоскость которых задаётся углом
  укладки, а у каждой пластинки --- около базиса.
\item \textbf{Количественно хуже.} Порог сдвинулся вниз ($F_{n45}=0.66$ уже при 125~МПа), плато
  RHF по снимку слишком высокое (0.96 при 200~МПа против 0.57 у Lepine), межзёренных меньше (0.55
  без нагрузки против 0.84--0.91 раньше), стопок мало (5--6 на поле).
\end{itemize}

\begin{table}[htbp]\centering\small
\caption{Модель с ореолами: $F_{n45}$ на снимке. Перекалибровка форы зёрен идёт.}
\begin{tabular}{lcccccc}\toprule
& 100 & 125 & 150 & 175 & 200~МПа & межзёренных при 0\\\midrule
Cinbiz & 0 & 0 & $\sim$0.05 & $\sim$0.95 & $\sim$1 & ---\\
ореолы, фора 10.9\,$^\circ$C & 0.17 & 0.66 & 0.93 & 0.93 & 0.99 & 0.55\\
ореолы, фора 13\,$^\circ$C (первый прогон) & 0.17 & 0.46 & 0.57 & 0.93 & 0.95 & ---\\
\bottomrule
\end{tabular}
\end{table}

\begin{warnbox}
Пока не проверено, не преувеличивает ли сетка 0.4~мкм автокатализ: клетка у самого кончика
получает очень сильное поле. Нужна проверка на сетке 0.2~мкм. Таблица ореолов --- только для
Zircaloy-4 ($\sigma_y=350$~МПа).
\end{warnbox}

\begin{exercise}
В модели с ореолами доля радиальных \emph{пластинок} при 250~МПа около 0.5--0.64, а RHF
\emph{по снимку} --- 0.99. Объясните, как так может быть, и какая из двух мер правильна для
сравнения с Lepine. Что из этого следует для отдельной пластинки и для стопки?
\end{exercise}

\section{Сверка численных инструментов}

Чтобы доверять выводам, проверены сами инструменты:
\begin{itemize}
\item упругое поле автомата против МКЭ --- совпадает до 0.5~МПа после исправления оси
  (гл.~\ref{ch:misfit});
\item упругопластический Фурье против МКЭ на стопках --- 1--21~МПа из 2700 (гл.~\ref{ch:plasticity});
\item сложение ореолов против честного расчёта --- ошибка $\sim60$~МПа на месте следующей
  пластинки, в три раза лучше потолка (гл.~\ref{ch:plasticity}).
\end{itemize}
